\section{Potential fields due to 3D sources}
\label{sec:potential_fields_3D_sources}


Let $\mathcal{V} \subset \mathbb{R}^{3}$ be a \textit{regular region} with boundary
$\partial \mathcal{V}$. 
Some important properties of $\mathcal{V}$ in this case is that it is an open and 
connected set of points and $\partial \mathcal{V}$ is a closed and compact surface 
free of double points \citep[][ p. 22]{freeden_gerhards2012}.
This definition is slightly different from that given by \citet[][ p. 100]{kellogg1967},
for example, which includes the boundary as part of the region.

Consider a set of 3D gravity or magnetic sources located in the subsurface,
over $\overline{\mathcal{V}} = \mathcal{V} \, \cup \, \partial \mathcal{V}$.
Consider also that the gravitational or magnetic induction field produced by the 3D sources 
at an observation point $\mathbf{r} = (x, y, z)^{\top}$ 
can be defined in terms of the following Newtonian potential:
\begin{equation}
	U(\mathbf{r}) = \kappa \iiint\limits_{\mathcal{V}}
	\sigma(\mathbf{r}') \: \frac{1}{\| \mathbf{r} - \mathbf{r}'\|}
	\: dx'dy'dz' \: ,
	\label{eq:scalar-potential-U}
\end{equation}
where $\kappa$ is a proportionality constant; 
$\tfrac{1}{\| \mathbf{r} - \mathbf{r}' \|}$ is the inverse distance function
between the observation point $\mathbf{r}$ and a point $\mathbf{r}' \in \mathcal{V}$;
and $\sigma(\mathbf{r}')$ is a scalar function defined at points 
$\mathbf{r}' = (x', y', z')^{\top} \in \overline{\mathcal{V}}$ as the density (contrast) 
or total-magnetization intensity within the sources.
To insure that the partial derivatives of first order of $U(\mathbf{r})$ can be computed
at points within the sources, $\sigma(\mathbf{r}')$ must be
a piecewise continuous, integrable and bounded function. 
Additionally, the partial derivatives of second order of $U(\mathbf{r})$ exist only at
points $\mathbf{r}' \in \overline{\mathcal{V}}$ where $\sigma(\mathbf{r}')$ satisfies some
Hölder condition \citep[][ p. 152]{kellogg1967}.
Hence, we impose these conditions on $\sigma(\mathbf{r}')$, so that the partial derivatives
of first and second orders of $U(\mathbf{r})$ exist at all points 
$\mathbf{r} \in \mathbb{R}^{3}$ in space, including those within the sources.

If we consider gravity sources, then $U(\mathbf{r})$ 
(equation \ref{eq:scalar-potential-U}) represents the gravitational potential
\citep[e.g.,][ p. 46]{blakely1996} and its partial derivatives of first order
\begin{equation}
	\partial_{\alpha} U(\mathbf{r}) = \kappa \iiint\limits_{\mathcal{V}}
	\sigma(\mathbf{r}') \: 
	\partial_{\alpha} \left( \frac{1}{\| \mathbf{r} - \mathbf{r}'\|} \right)
	\: dx'dy'dz' \: , \quad \alpha = x, y, z \: ,
	\label{eq:1st-deriv-U}
\end{equation}
represent the Cartesian components of the gravitational attraction 
\citep[e.g.,][ p. 47]{blakely1996} due to the 3D sources.
In geophysics, the vertical component of the gravitational attraction defined by
$\partial_{z} U(\mathbf{r})$ (equation \ref{eq:1st-deriv-U}) is commonly used to approximate
\textit{gravity disturbance} data \cite[e.g.,][ p. 93]{hofmann-wellenhof-moritz2005}.

The partial derivatives of second order of $U(\mathbf{r})$ 
(equation \ref{eq:scalar-potential-U}) are given by:
\begin{equation}
	\partial_{\alpha\beta} U(\mathbf{r}) = \kappa \iiint\limits_{\mathcal{V}}
	\sigma(\mathbf{r}') \: 
	\partial_{\alpha\beta} \left( \frac{1}{\| \mathbf{r} - \mathbf{r}'\|} \right)
	\: dx'dy'dz' \: , \quad \alpha, \beta = x, y, z \: ,
	\label{eq:2nd-deriv-U}
\end{equation}
and represent the elements of the \textit{gravitational tensor} 
\citep[e.g.,][]{pedersen_rasmussen1990} produced by the 3D sources. 
Note that we use the notation 
$\partial_{\alpha} U(\mathbf{r}) \equiv 
\tfrac{\partial U(\mathbf{r})}{\partial \alpha}$ and
$\partial_{\alpha} U(\mathbf{r}) \equiv 
\tfrac{\partial^{2} U(\mathbf{r})}{\partial \alpha \, \partial \beta}$
to denote, respectively, partial derivatives of first and second order throughout the
manuscript.

To define the magnetic fields due to the 3D sources, we limit ourselves to the case
in which they have a constant total-magnetization direction defined by an unit vector
$\hat{\mathbf{h}}$. In this case, the magnetic scalar potential is given by the
negative of the following directional derivative:
\begin{equation}
	\partial_{h} U(\mathbf{r}) = \nabla U(\mathbf{r})^{\top} \hat{\mathbf{h}} \: ,
	\label{eq:oblique-1st-deriv-U}
\end{equation}
where 
$\partial_{h} U(\mathbf{r}) \equiv 
\tfrac{\partial U(\mathbf{r})}{\partial \hat{\mathbf{h}}}$ and 
$\nabla U(\mathbf{r})$ represents the gradient of $U(\mathbf{r})$ with respect to the
coordinates of the observation points $\mathbf{r}$.

The magnetic induction field along an arbitrary constant direction represented
by an unit vector $\hat{\mathbf{t}}$ can be defined as follows:
\begin{equation}
	\partial_{th} U(\mathbf{r}) = 
	\hat{\mathbf{t}}^{\top} \mathbf{H}_{U}(\mathbf{r}) \, \hat{\mathbf{h}} \: ,
	\label{eq:oblique-2nd-deriv-U}
\end{equation}
where 
$\partial_{th} U(\mathbf{r}) \equiv 
\tfrac{\partial^{2} U(\mathbf{r})}{\partial \hat{\mathbf{h}} \, \partial \hat{\mathbf{t}}}$
and
\begin{equation}
	\mathbf{H}_{U}(\mathbf{r}) = 
	\begin{bmatrix}
	\partial_{xx} U(\mathbf{r}) & 
	\partial_{xy} U(\mathbf{r}) & 
	\partial_{xz} U(\mathbf{r}) \\
	\partial_{xy} U(\mathbf{r}) & 
	\partial_{yy} U(\mathbf{r}) & 
	\partial_{yz} U(\mathbf{r}) \\
	\partial_{xz} U(\mathbf{r}) & 
	\partial_{yz} U(\mathbf{r}) & 
	\partial_{zz} U(\mathbf{r})
	\end{bmatrix}
	\label{eq:H-U}
\end{equation}
is the Hessian matrix of $U(\mathbf{r})$ (equation \ref{eq:scalar-potential-U}), 
with elements defined by its second derivatives (equation \ref{eq:2nd-deriv-U}).
Equation \ref{eq:oblique-2nd-deriv-U} is commonly used to approximate the 
\textit{total-field anomaly} produced by the 3D sources for the
case in which $\hat{\mathbf{t}}$ defines the constant direction of the main geomagnetic
field at the study area \cite[e.g., ][ p. 179]{blakely1996}.

The $x-$, $y-$, and $z-$components of the magnetic induction field produced by the 3D
sources can be easily obtained from equation \ref{eq:oblique-2nd-deriv-U}
by setting the unit vector $\hat{\mathbf{t}}$ equal to,
respectively, $(1, 0, 0)^{\top}$, $(0, 1, 0)^{\top}$, and $(0, 0, 1)^{\top}$.
Finally, the particular case in which both unit vectors $\hat{\mathbf{h}}$ and 
$\hat{\mathbf{t}}$ are set as $(0, 0, 1)^{\top}$ in equation \ref{eq:oblique-2nd-deriv-U}
leads to a special magnetic anomaly defined by the partial derivative of second order 
$\partial_{zz}U(\mathbf{r})$ (equation \ref{eq:2nd-deriv-U}). 
This special anomaly represents the vertical component of 
the magnetic induction field that would be produced by the 3D sources if they had a 
purely induced and vertical magnetization. Differently from generic total-field anomalies 
\cite[e.g.,][ p. 178]{blakely1996}, which have a ``skewed'' pattern 
oscillating from negative to positive values and do not have their maxima placed right 
above the magnetic sources, this special magnetic anomaly is easier to interpret 
because it is predominantly positive right above the sources and
decays to zero towards their boundaries. \citet{baranov1957} was the first
to propose a method, called \textit{reduction to the pole} (RTP), to transform generic
total-field anomalies into this special anomaly, conveniently called here 
\textit{RTP anomaly}.




